\chapter{Zusammenfassung und zukünftige Arbeiten}

\section{Zusammenfassung der Forschung und Beiträge}
Diese Arbeit hat den architektonischen Entwurf, die Implementierung und die
Validierung eines paravirtualisierten VirtIO-GPU-Treibers für den
hhuOS-Forschungskernel detailliert beschrieben. Der primäre Beitrag dieser
Arbeit ist die erfolgreiche Verbindung eines leistungsfähigen
Grafikprotokolls mit einer Mikrokernel-Umgebung -- ausgehend von einer
statischen Testausgabe bis hin zu einem vollständigen dynamischen
Laufzeit-Rendering-Stack .

Durch diese Forschung wurde gezeigt, dass die Komplexität moderner Grafik durch
Paravirtualisierung beherrscht werden kann, ohne die Kernprinzipien der
Sicherheit und Isolierung zu opfern. Die Implementierung des
\texttt{virtio\_gpu}-Treibers hat hhuOS in ein System erweitert, das einen
Durchsatz von 564 FPS erreicht und damit eine 7,4-fache Verbesserung gegenüber
Legacy-Methoden erzielt. Die dynamische Erweiterung -- bestehend aus dem
Per-Frame-Renderer (\texttt{virtio\_gpu\_demo\_render\_frame}), der
Dirty-Rect-Präsentation (\texttt{present\_demo\_regions}), einem dedizierten
Animations-Thread und einem interaktiven Laufzeit-Moduswechsel über die Taste
\texttt{C} -- erhöht den Reifegrad des Treibers erheblich und validiert die
gesamte Rendering-Pipeline unter realistischen dynamischen Bedingungen.

\section{Technische Reflexion: Paravirtualisierung vs. Emulation}
Die dokumentierten Ergebnisse liefern empirische Belege dafür, dass das
„Trap-and-Emulate“-Modell für moderne grafische Arbeitslasten nicht mehr
tragfähig ist. Die Implementierung des linearen Framebuffers
(Legacy Linear Framebuffer, LFB) ist zwar einfach umzusetzen, verursacht jedoch
aufgrund ständiger VM-Exits eine erhebliche Leistungseinbuße.

Die VirtIO-GPU-Implementierung zeigt, dass die
„Transfer-then-Flush“-Semantik überlegen ist. Die zusätzliche
Dirty-Rect-Präsentation verstärkt diesen Effekt, indem der Hypervisor nur für
tatsächlich veränderte Bildschirmregionen aktiviert wird. Die Validierung der
spezifischen Befehle und der Queue-Strukturen zeigt, dass die Reduzierung der
„semantischen Lücke“ ein wesentlicher Schlüssel zur paravirtualisierten
Effizienz ist.

\section{Architektonische Auswirkungen auf hhuOS}
Aus softwaretechnischer Perspektive hat sich die modulare Struktur des Treibers
als äußerst effektiv erwiesen. Die Implementierung ist in mehrere Submodule
unter \texttt{src/device/graphic/virtio/} gegliedert. Diese Struktur trennt
Discovery, Initialisierung, Virtqueue-Transport, Ressourcenverwaltung und
Laufzeitpfade klar voneinander. Dadurch bleibt der Kernel übersichtlich und
prüfbar, während gleichzeitig komplexe Hardware strukturiert unterstützt werden
kann.

Die in den Tabellen~5.1 bis~5.4 zusammengefassten Strukturen und Funktionen
dienen als Blaupause für zukünftige VirtIO-basierte Treiber in hhuOS, wie
beispielsweise \texttt{virtio-net} oder \texttt{virtio-blk}, die den gleichen
hier etablierten Virtqueue-basierten Kommunikationsmustern folgen können.

Besonders hervorzuheben ist die Erweiterung durch den dedizierten
Animations-Thread: Sie zeigt, dass hhuOS in der Lage ist, nebenläufige
Hintergrundaufgaben mit direktem Hardwarezugriff zu verwalten -- ein
wesentlicher Schritt in Richtung eines vollständigen Grafiksubsystems.

\section{Einschränkungen und offene Herausforderungen}
Trotz des Erfolgs des aktuellen Treibers verbleiben einige Einschränkungen. Die
aktuelle Implementierung ist auf 2D-Operationen beschränkt. Obwohl das
2D-Protokoll stabil und leistungsfähig ist, sind moderne Desktop-Umgebungen
zunehmend auf 3D-Shader und Vertex-Verarbeitung angewiesen.

Darüber hinaus arbeitet der Treiber derzeit auf Ebene der Control Queue
synchron: Die CPU muss den Used Ring abfragen und auf die Antwort des Hosts
warten. Dies gewährleistet zwar Datenkonsistenz, kostet jedoch zusätzliche
CPU-Zeit. Der Übergang zu einem interruptgesteuerten oder weiter
asynchronisierten Modell würde die Systemreaktionsfähigkeit weiter verbessern.

\section{Zukünftige Arbeiten: In Richtung 3D und darüber hinaus}
Der logische nächste Schritt für diese Forschung ist die Integration von
\textbf{VirGL}. VirGL ist eine Erweiterung des VirtIO-GPU-Protokolls, die es
dem Gast ermöglicht, OpenGL-Befehle direkt an den Host zu senden. Dies würde
hhuOS in die Lage versetzen, hardwarebeschleunigte 3D-Anwendungen
auszuführen.

Darüber hinaus sollten zukünftige Arbeiten die Nutzung von
\textbf{IOMMU (Input-Output Memory Management Unit)}-Unterstützung innerhalb
des hhuOS-Kernels untersuchen. Dies würde eine zusätzliche Sicherheitsebene
hinzufügen und sicherstellen, dass selbst bei einer Kompromittierung des
Grafiktreibers kein Zugriff auf Speicherbereiche außerhalb der vorgesehenen
DMA-Puffer möglich ist.

\section{Abschließende Bemerkungen}
Der Übergang von textbasierten zu grafischen Oberflächen ist ein Meilenstein
für jedes Betriebssystem. Durch die Implementierung eines
VirtIO-GPU-Treibers -- vom statischen Schachbrettmuster bis zur vollständig
animierten und interaktiv umschaltbaren Rendering-Pipeline -- wurde hhuOS nicht
nur um eine moderne Grafikausgabe erweitert, sondern zugleich die Fähigkeit des
Kernels validiert, Gerätekommunikation mit hohem Durchsatz und niedriger
Latenz zu bewältigen. Da die Virtualisierung weiterhin sowohl im Rechenzentrum
als auch auf dem Desktop eine zentrale Rolle spielt, werden die in dieser
Arbeit dokumentierten Techniken für den Entwurf sicherer, effizienter und
moderner Betriebssysteme weiterhin von grundlegender Bedeutung bleiben.